\section{Introduction}\label{sec:intro}

循环融合能够避免重复遍历和物化中间结果~\citep{halide,tvm}。
但是，当后续循环读取较早归约的最终结果时，通过交错执行迭代来融合这些循环会违反这一依赖。
这类依赖出现在机器学习和科学计算的张量算子中。

% Shared unchanged by both prose editions.
\begin{wrapfigure}{r}{0.33\textwidth}
\vspace{-6pt}
\centering
\captionsetup{skip=3pt}
\begin{tikzpicture}[x=1cm,y=1cm,
 instance/.style={draw=origblue,rounded corners=2pt,fill=boxgray,inner sep=2pt,font=\scriptsize},
 carry/.style={->,draw=rulegray},
 dependency/.style={->,draw=staleorange,semithick}]
\node[instance] (a) at (0,0) {$I_1(0)$};
\node[instance] (b) at (1.15,0) {$I_1(1)$};
\node[font=\scriptsize,inner sep=1pt] (gap1) at (2.10,0) {$\cdots$};
\node[instance] (c) at (3.15,0) {$I_1(N-1)$};
\node[instance] (d) at (0,-0.95) {$I_2(0)$};
\node[instance] (e) at (1.15,-0.95) {$I_2(1)$};
\node[font=\scriptsize,inner sep=1pt] (gap2) at (2.10,-0.95) {$\cdots$};
\node[instance] (f) at (3.15,-0.95) {$I_2(N-1)$};
\node[font=\scriptsize,anchor=south west,inner sep=0pt]
  at ([yshift=2pt]a.north west) {Loop 1};
\node[font=\scriptsize,anchor=south west,inner sep=0pt]
  at ([yshift=2pt]d.north west) {Loop 2};
\foreach \s/\t in {a/b,b/gap1,gap1/c,d/e,e/gap2,gap2/f} \draw[carry] (\s) -- (\t);
\foreach \t in {d,e,f} \draw[dependency] (c.south) -- (\t.north);
\end{tikzpicture}
\caption{Schematic of dependent reductions.}
\label{fig:dependent-loops}
\vspace{-6pt}
\end{wrapfigure}

在\emph{依赖归约}中，归约体使用由同一轴上较早归约的最终结果确定的值。
这样的值是它的\emph{上下文}：它沿归约轴保持不变，但只有在较早计算完成后才可用。
\cref{fig:dependent-loops} 中，$I_i(j)$ 表示循环 $i$ 的第 $j$ 次迭代，橙色箭头表示循环之间的依赖。

因此，要将这些归约融合为单遍计算，就需要在最终上下文尚不可用时，先为后续归约积累信息。
为一个输入分块保留的聚合量是\emph{状态分量}，简称\emph{状态}；保留的量共同组成\emph{状态元组}。

\emph{合并}（join）将两个不相交分块的元组组合为其并集对应的元组。
对于使用局部上下文计算的结果，这一操作必须考虑每个分块的上下文变为并集上下文所带来的影响。
由此产生两个相互关联的问题：哪些状态能够保留足够的信息，以及如何推导正确更新这些状态的合并操作。

开发这类算法需要专门经验。
在线归一化、矩累积和分块注意力都需要识别保留量、推导更新，并实现和调优内核~\citep{milakov,chan,pebay,flashattention}。
近期的自动融合方法能够处理平移或缩放局部结果等代数修正~\citep{neptune,redfuser}。
当修正依赖于原有归约结果未保留的聚合信息时，这些规则便不再适用。

\textbf{Our Work.}
本文提出 \sys，\footnote{Greek \emph{hapax}, ``once''.} 通过扩充源归约结果并推导其合并操作，融合依赖张量归约。
我们借鉴为分治并行开发的方法：Farzan 和 Nicolet 综合辅助累加量与合并操作~\citep{farzan17,farzan21}，AutoLifter 则为分治计算综合辅助值和组合函数~\citep{autolifter}。
这些方法引入额外信息，使各分块独立计算的结果能够组合。
\sys 将这一思想用于融合：\emph{辅助状态}补足源结果的上下文修正所需的信息。

\par\medskip
\noindent\begin{minipage}{\linewidth}
\begin{minipage}[t]{0.50\linewidth}
\vspace{0pt}
各分块使用自身的上下文计算状态元组，合并操作再修正并组合这些元组（\cref{fig:dc-fusion}）。
尽管较大分块的元组可能需要多次遍历才能得到，单元素分块的元组只依赖一个输入元素。
在这一构造中，{\bfseries 融合循环}可以视为在一次遍历中依次合并单元素分块的状态元组。
\end{minipage}\hfill
% Shared English vector figure; both prose editions include this source.
% The introducing prose occupies the other half of the enclosing minipage.
\begin{paperdiagram}
\captionsetup{type=figure,skip=4pt}
\centering
\begin{tikzpicture}[x=1cm,y=1cm,>={Stealth[length=1.2mm,width=0.9mm]},
  label/.style={font=\footnotesize,align=center},
  instance/.style={draw=origblue,rounded corners=2pt,fill=boxgray,
    inner sep=2pt,minimum height=0.32cm,font=\footnotesize},
  continuation/.style={font=\footnotesize,inner sep=1pt},
  part/.style={draw=origblue,rounded corners=4pt},
  joinop/.style={draw=fixgreen,rounded corners=2pt,fill=fixgreen!8,
    inner xsep=6pt,inner ysep=3pt,font=\footnotesize\bfseries},
  carry/.style={->,draw=rulegray},
  dependency/.style={->,draw=staleorange,semithick},
  combine/.style={->,draw=fixgreen,semithick}]

\node[label] at (1.50,0) {Partition $T$};
\node[label] at (4.90,0) {Partition $S$};
\draw[part] (0,-0.22) rectangle (3.00,-2.10);
\draw[part] (3.40,-0.22) rectangle (6.40,-2.10);

% Symbolic endpoints stand for the first and last iterations of each part.
\node[instance] (t1a) at (0.55,-0.55) {$I_1(a)$};
\node[continuation] (t1dots) at (1.50,-0.55) {$\cdots$};
\node[instance] (t1b) at (2.45,-0.55) {$I_1(b)$};
\node[instance] (t2a) at (0.55,-1.50) {$I_2(a)$};
\node[continuation] (t2dots) at (1.50,-1.50) {$\cdots$};
\node[instance] (t2b) at (2.45,-1.50) {$I_2(b)$};
\node[instance] (s1c) at (3.95,-0.55) {$I_1(c)$};
\node[continuation] (s1dots) at (4.90,-0.55) {$\cdots$};
\node[instance] (s1d) at (5.85,-0.55) {$I_1(d)$};
\node[instance] (s2c) at (3.95,-1.50) {$I_2(c)$};
\node[continuation] (s2dots) at (4.90,-1.50) {$\cdots$};
\node[instance] (s2d) at (5.85,-1.50) {$I_2(d)$};
\foreach \a/\dots/\b in {t1a/t1dots/t1b,t2a/t2dots/t2b,s1c/s1dots/s1d,s2c/s2dots/s2d} {
  \draw[carry] (\a.east) -- (\dots.west);
  \draw[carry] (\dots.east) -- (\b.west);
}
% Each part's completed first reduction supplies its second reduction.
\foreach \a/\b in {t1b/t2a,t1b/t2b,s1d/s2c,s1d/s2d}
  \draw[dependency] (\a.south) -- (\b.north);
\node[label,text=auxpurple] at (1.50,-1.90) {Supplemental states};
\node[label,text=auxpurple] at (4.90,-1.90) {Supplemental states};

\node[joinop] (combined) at (3.20,-2.70) {Join};
\draw[combine] (1.50,-2.10) -- (combined.north west);
\draw[combine] (4.90,-2.10) -- (combined.north east);
\end{tikzpicture}
\caption{Partitioned computation.}
\label{fig:dc-fusion}
\end{paperdiagram}

\end{minipage}\par\medskip

为确定哪些信息能够支持这样的合并，\sys 通过代数分析将归约体与其上下文分离。
所得\emph{参考状态}是归约体不读取这些上下文的归约，因而\emph{不依赖上下文}。
它们能够独立合并，并确定所需结果，所以分析成功时便得到一种单遍构造。
其数量给出不随归约长度变化的状态数充分上界。

参考构造提供了一种有效的单遍表示，但执行其展开算术可能破坏源程序的中心化或缩放，引入消去或溢出~\citep{chan,blanchard,boldo_fp}。
这些数值方面的考虑促使我们以源归约在局部上下文下计算的结果作为\emph{维护状态元组}的起点；该元组包含融合程序实际累积的量。
保留这些结果，使源程序中经过中心化或缩放的量能够继续用于上下文修正。
当源结果未保留这些修正所需的足够信息时，辅助状态补足缺失的信息。

\sys 综合辅助状态，使其与源归约结果共同组成能够与参考元组相互转换的维护元组。
对于每个候选维护元组，综合过程构造并验证相应的转换表达式。
通过恢复各分块的参考元组、执行已知的参考合并，再重构并集上的维护元组，即可得到正确的合并操作。
因此，综合只需寻找满足这一转换要求的辅助状态，无需独立搜索合并表达式。

为限制辅助状态的搜索空间，我们根据源归约结果对上下文变化的响应构造候选表达式。
符号化地平移或缩放局部上下文，会引出修正结果所需的聚合系数。
这些系数为语法引导综合提供候选表达式，将搜索限定在对更新有明确作用的量上。
固定的源分量和这些候选共同限制了综合器需要发现的内容。
同样的修正还提供执行公式，在调整局部结果时尽可能复用源程序的中心化和缩放。
以通过恢复得到的合并操作为依据验证这些公式，将数值计算形式与代数正确性论证联系起来。

在实数语义下，所得合并操作计算组合后输入的维护状态元组。
由于合并也适用于任意不相交的非空输入部分，所得归约器允许各分块独立计算后通过归约树组合。
减少遍历为降低数据搬运开销创造了机会；实际性能还取决于状态大小、合并成本和调度。

我们将 \sys 实现为 TVM 张量表达式（TE）的重写~\citep{tvm}。
一个已记录的实现快照在机器学习和科学计算中选取的 \nShowcase 个算子中，为 \nShowFused 个推导出单遍归约器。
本文报告覆盖率、发现的状态和推导开销，并介绍实现中用于比较生成程序与源程序精度的 MPFR 检验方法。
GPU 调度和性能评估仍是后续工作。

\noindent 本文的贡献如下：
\begin{itemize}
\item \textbf{Fusion of dependent-reduction patterns.}
上下文分离识别出一类具有单遍构造的依赖归约链，其中包括消费者归约体超出近期自动融合方法所述代数条件的情形，并界定这一构造所需状态数的充分上界。
\item \textbf{Reference-guided state synthesis and join derivation.}
完整恢复参考表示将综合缩减为状态扩充，使独立搜索合并表达式变得不再必要。
\item \textbf{Source-guided search and numerical formulation.}
源结果及其上下文修正限制辅助状态搜索，同时提供尽可能复用源程序中心化和缩放的执行公式。
\item \textbf{Implementation and evaluation.}
TVM 实现，以及 \nShowcase 个算子上已记录的覆盖率、发现的状态和推导开销结果。
\end{itemize}
